\section{A graded obstruction to coordinateness}
\label{sec:filtration}

We now turn to the assertion that $p$ is not a coordinate of $A$.
The first step is to pass to a weighted degeneration of $A$.
If $p$ were a coordinate, differentiation in another coordinate would
then produce a nonzero locally nilpotent derivation of a simpler
quadric ring, with a nonzero invariant in an explicit ideal.
The rest of the proof will exclude that possibility.
The use of leading derivations follows the filtration methods of
\cite[Sections~3--5]{KalimanMakarLimanov2007}; the particular quadric construction is
adapted from \cite[Sections~3--5]{Cancellation2026}.

\subsection{The associated graded ring}

Give the generators of $P=k[p,s,u,F,J]$ the weights
\begin{equation}
\label{gr:weights}
\begin{array}{c|rrrrr}
  &p&s&u&F&J\\ \hline
 \text{weight}&-1&0&0&1&1.
\end{array}
\end{equation}
For $d\in\mathbb Z$, let $\mathcal F_dP$ be the span of the monomials
of weight at most $d$, and give $A=P/(H)$ the quotient filtration
$\mathcal F_dA=\operatorname{im}(\mathcal F_dP\to A)$.
These are increasing multiplicative filtrations, indexed by all
integers.  Their negative indices are essential: $p$ has weight $-1$.

\begin{proposition}
\label{prop:graded}
The filtration on $A$ is exhaustive and separated.  Its associated
graded ring is the domain
\begin{equation}
\label{gr:presentation}
 G:=\gr A
 \simeq k[p,s,u,F,J]/(H_*),
 \qquad
 H_*=x^2F-(1+2sx)J-pJ^2,
 \quad x=s^2-u^2+pF,
\end{equation}
with the grading induced by \eqref{gr:weights}.
The filtration degrees of $p,s,u,F,J$ in $A$ are respectively
$-1,0,0,1,1$.
\end{proposition}

\begin{proof}
Exhaustiveness follows because every element of $A$ has a polynomial
representative.  For $r>0$, every monomial of weight at most $-r$ is
divisible by $p^r$, so
\[
 \mathcal F_{-r}A\subseteq p^rA.
\]
By Proposition~\ref{prop:polynomial}, $A$ is a polynomial ring, and
Equation~\eqref{con:candidate} shows that $p$ is a nonzero nonunit.
Therefore $\bigcap_{r>0}p^rA=0$: in a unique
factorization domain, a nonzero element cannot be divisible by
arbitrarily high powers of an irreducible factor of $p$.
This proves separation.  In particular, every nonzero $a\in A$ has a
well-defined degree
\[
 \deg a=\min\{d\in\mathbb Z:a\in\mathcal F_dA\}
\]
and a nonzero initial form in $\mathcal F_dA/\mathcal F_{d-1}A$.

The element $x$ is homogeneous of weight zero in $P$, and
$H=H_*-u$ has highest weight part $H_*$, of weight one.
For every nonzero $q\in P$, the highest weight part of $qH$ is
$(\operatorname{in}q)H_*$: the polynomial ring is a graded domain,
so the product of the highest parts does not vanish.
Consequently the ideal of initial relations is exactly $(H_*)$.
For completeness, the natural graded map $P\to\gr A$ sends a homogeneous
polynomial of weight $d$ to its class in
$\mathcal F_dA/\mathcal F_{d-1}A$, which may at this stage be zero.
This map is surjective.
A homogeneous polynomial of weight $d$ lies in its kernel exactly
when its image in $A$ has a representative of weight at most $d-1$;
subtracting that representative gives a relation with the original
polynomial as its highest part.  Thus the kernel is the ideal of
initial relations, proving \eqref{gr:presentation}.

To prove that this quotient is a domain, use the polynomial
coordinates $x,y,z,u$ of $P[p^{-1}]$ from the preceding section.
The identity $xy-z(z+1)=p(H+u)$ becomes
\[
 pH_*=xy-z(z+1).
\]
The polynomial on the right is irreducible over
$k[p,p^{-1},x,z,u]$: as a polynomial in $y$, it is primitive,
because $x$ and $z(z+1)$ are coprime, and it has degree one over
the fraction field.  Thus $H_*$ is irreducible after inverting $p$.
It is not divisible by $p$ before localization, since its reduction
modulo $p$ is
\[
 (s^2-u^2)^2F-\bigl(1+2s(s^2-u^2)\bigr)J\ne0.
\]
Any factorization of $H_*$ in $P$ has a factor that becomes a unit after
inverting $p$, hence is a scalar times a power of $p$.
Since $p\nmid H_*$, this factor is already a unit in $P$.
Thus $H_*$ is irreducible in $P$, and hence prime by unique factorization.
Finally, $H_*$ is not associated to any
of the five variables, so their classes in $G$ are nonzero.
These classes are the initial forms of the corresponding elements
of $A$, proving the asserted degrees.
\end{proof}

\subsection{A quadric and its distinguished ideal}

From now on write $\tau$ for the class of $p$ in $G$, to distinguish
the graded ring from $A$.  Consider the quadric ring
\begin{equation}
\label{gr:quadric}
 B=k[x,y,z,u]/\bigl(xy-z(z+1)\bigr).
\end{equation}
The formulas
\[
 x=s^2-u^2+\tau F,\qquad
 y=s+x(x+u^2),\qquad z=sx+\tau J
\]
define elements of degree zero in $G$.  They satisfy the relation in
\eqref{gr:quadric}; after inverting $\tau$, the inverse coordinate
formulas from the preceding section give an isomorphism
\begin{equation}
\label{gr:localization}
 G[\tau^{-1}]=B[\tau,\tau^{-1}],
 \qquad \deg B=0,\quad \deg\tau=-1.
\end{equation}
In particular, the induced map $B\to G$ is injective: its composition
with localization is the inclusion of $B$ in this Laurent polynomial
ring.  We henceforth regard $B$ as a subring of $G$.

Define elements and an ideal of $B$ by
\begin{equation}
\label{gr:ideal}
 \begin{aligned}
 s_B&=y-x(x+u^2),\\
 f_0&=x-s_B^2+u^2,\qquad
 g_0=z-s_Bx,\qquad I=(f_0,g_0).
 \end{aligned}
\end{equation}
In \eqref{gr:localization} we have $s=s_B$, $F=f_0/\tau$, and
$J=g_0/\tau$.  These identities give the more precise description
\begin{equation}
\label{gr:modification}
 G=B[\tau,I/\tau]
   =B[\tau,f_0/\tau,g_0/\tau]
   \subset B[\tau,\tau^{-1}].
\end{equation}
This is the algebraic construction of an affine modification, with
divisor $\tau=0$ and center $(\tau,I)$ in $B[\tau]$; see
\cite[Definition~1.1 and Proposition~1.1(b)]{KalimanZaidenberg1999}.
We will only need its explicit
homogeneous pieces.

\begin{lemma}
\label{gr:positive-piece}
For each integer $\ell>0$, the degree-$\ell$ piece of $G$ is
\[
 G_\ell=\tau^{-\ell}I^\ell.
\]
\end{lemma}

\begin{proof}
A monomial in the generators in \eqref{gr:modification}, with
coefficient in $B$, has the form
\[
 b\,\tau^r(f_0/\tau)^i(g_0/\tau)^j,
 \qquad r,i,j\ge0.
\]
It has degree $\ell$ exactly when $i+j-r=\ell$.
In that case it equals $\tau^{-\ell}b f_0^i g_0^j$, whose numerator
belongs to $I^{\ell+r}\subseteq I^\ell$.
Since the presentation is homogeneous, such monomials span $G_\ell$.
Conversely, every element of $I^\ell$ is a $B$-linear combination of
the products $f_0^i g_0^j$ with $i+j=\ell$, and dividing these
products by $\tau^\ell$ gives elements of $G_\ell$.
\end{proof}

\subsection{What a coordinate would force}

The ideal $I$ records the positive part of the grading.  We now show
why a hypothetical coordinate system containing $p$ supplies a nonzero
locally nilpotent derivation of $B$ fixing a nonzero element of $I$.

\begin{proposition}
\label{prop:coordinate-obstruction}
If $p$ is a coordinate of $A$, then there exist a nonzero locally
nilpotent derivation $E$ of the ring $B$ in \eqref{gr:quadric} and
an element $0\ne h\in I$ such that $E(h)=0$.
\end{proposition}

\begin{proof}
Suppose $A=k[p,q_1,q_2,q_3]$ for a polynomial coordinate system.
At least one $q_i$ has strictly positive filtration degree.
Indeed, if all three had degree at most zero, then so would every
polynomial in $p,q_1,q_2,q_3$, contrary to $\deg F=1$.
Relabel so that $\deg q_1>0$, and take
\[
 D=\frac{\partial}{\partial q_2}.
\]
This is a nonzero locally nilpotent derivation of $A$ satisfying
$D(p)=D(q_1)=0$.

We give the details of its passage to $G$, because the filtration
has both positive and negative indices.
Write $\omega(t)$ for the assigned weight of a generator
$t\in\{p,s,u,F,J\}$, and put
\begin{equation}
\label{gr:shift}
 m=\max_{D(t)\ne0}\bigl(\deg D(t)-\omega(t)\bigr).
\end{equation}
This maximum is defined: a nonzero derivation cannot kill all
algebra generators.  The product rule shows that a monomial of
weight at most $d$ has derivative in $\mathcal F_{d+m}A$.
Every element of $\mathcal F_dA$ has a representative that is a sum
of such monomials, by the definition of the quotient filtration.
Thus
\[
 D(\mathcal F_dA)\subseteq\mathcal F_{d+m}A
 \qquad(d\in\mathbb Z).
\]
It follows that $D$ induces a homogeneous derivation $D_0$ of $G$
of degree $m$.
For $a\in\mathcal F_dA$, write $[a]_d$ for its class in
$\mathcal F_dA/\mathcal F_{d-1}A$.  Explicitly,
\begin{equation}
\label{gr:iterate}
 D_0([a]_d)=[D(a)]_{d+m},\qquad
 D_0^r([a]_d)=[D^r(a)]_{d+rm}\quad(r\ge0).
\end{equation}
The first formula is well-defined because $D$ also sends
$\mathcal F_{d-1}A$ into $\mathcal F_{d+m-1}A$; the second follows
by induction.  It proves that $D_0$ is locally nilpotent, first on
homogeneous elements and then on finite sums of them.
Moreover, equality is attained in \eqref{gr:shift} for some
generator $t$, whose initial form has degree $\omega(t)$ by
Proposition~\ref{prop:graded}.  Its image under $D_0$ is the nonzero
initial form of $D(t)$, so $D_0\ne0$.
Finally, \eqref{gr:iterate} shows that $D_0$ kills both $\tau$ and
the nonzero initial form $\eta$ of $q_1$.  The latter is homogeneous
of degree $\ell=\deg q_1>0$.

Because $D_0(\tau)=0$, extending $D_0$ to
$G[\tau^{-1}]=B[\tau,\tau^{-1}]$ preserves local nilpotence.
Homogeneity implies that there is a derivation $E:B\to B$ with
\begin{equation}
\label{gr:descent}
 D_0(b)=\tau^{-m}E(b)\qquad(b\in B).
\end{equation}
Indeed, every degree-$m$ element of the Laurent ring is uniquely
$\tau^{-m}$ times an element of $B$.
The derivation $E$ is nonzero: otherwise $D_0$ would kill both $B$
and $\tau$, hence the whole Laurent ring and its subring $G$.
It is locally nilpotent because
\[
 E^r(b)=\tau^{rm}D_0^r(b)
 \qquad(b\in B, r\ge0).
\]
Here, as in \eqref{gr:iterate}, the formula follows by induction,
using that $D_0$ fixes $\tau$.

By Lemma~\ref{gr:positive-piece}, write
$\eta=\tau^{-\ell}h$ with $0\ne h\in I^\ell$.
Equations \eqref{gr:descent} and $D_0(\eta)=0$ give
\[
 0=D_0(\eta)=\tau^{-m-\ell}E(h).
\]
Thus $E(h)=0$, and $I^\ell\subseteq I$ because $\ell>0$.
\end{proof}

Proposition~\ref{prop:coordinate-obstruction} has removed the
original coordinate system from the problem.  It remains to show
that no nonzero locally nilpotent derivation of $B$ can have a
nonzero invariant in $I$.  We will make this ideal principal after
pulling it back to the complement of the zero section in a line bundle,
where a second grading gives the contradiction.
